\expandafter\def\csname HistoricalNoteText376\endcsname{Die erste öffentliche Erwähnung dieser Versuche findet man in den Gött. gel. Anz. 1834, Aug. 9. Vergl. \textsc{Schumachers} Jahrbuch für 1836, [S. 339 d. B.].}
\expandafter\def\csname HistoricalNoteText376b\endcsname{Gött. gel. Anz. 1835, März 7. \textsc{Schumachers} Jahrbuch für 1836, [S. 341 d. B.]}
\par
\SourceLeaf{372}Durch die Darstellung der Änderungen der Declination und der Intensität in zwei besondern Curven erhält man übrigens von dem Hergange der Störungen ein lange nicht so anschauliches Bild, wie durch ihre Vereinigung in eine einzige. Auf was es dabei ankommt, übersieht man am klarsten auf folgende Art. Eine vollständige Vorstellung der erdmagnetischen Kraft (nemlich des horizontalen Theils, wie immer stillschweigend verstanden wird) in jedem Augenblick kann man durch Eine gerade Linie geben, deren Länge der Intensität proportional ist, und die mit einer festen geraden Linie einen der Declination gleichen Winkel macht. Zur Darstellung der Kraft in mehrern auf einander folgenden Zeitpunkten lässt man den Anfangspunkt der verschiedenen geraden Linien unverändert, so dass die Endpunkte allein in Betracht kommen, die dann mit den entsprechenden Zeiten bezeichnet, und durch eine Linie vereinigt werden können. Die geraden Radien selbst werden gar nicht mitgezeichnet, und selbst der gemeinschaftliche Anfangspunkt wird bei einem nur einigermassen schicklichen Maassstab für die Darstellung immer weit ausserhalb der Zeichnung liegen. Diese Behandlung führt uns zugleich auf einen neuen Gesichtspunkt, aus welchem wir solche Veränderungen der beiden magnetischen Elemente betrachten können. Sie sind in der That nur die beiden horizontalen Componenten derjenigen vergleichungsweise immer sehr kleinen störenden Kraft, welcher in jedem Augenblick die mittlere erdmagnetische Kraft unterworfen ist, indem nemlich jene in zwei Richtungen, die eine \textit{im} magnetischen Meridian, die andere senkrecht gegen denselben zerlegt wird. Die zweite Componente wird unmittelbar durch das Magnetometer, die erste durch den neuen Apparat gegeben, wobei nur beide vor der Zeichnung auf ein gemeinschaftliches Maass zurückgeführt werden müssen.

Nur ein Umstand bei der Anwendung dieser an sich so anschaulichen Darstellungsart muss hier noch berührt werden, nemlich dass es nicht gut angeht, den Verlauf für einen ganzen Tag in Einer Zeichnung ohne Verwirrung darzustellen, wenn häufige schnell wechselnde Störungen vorkommen, da in diesem Fall die Curve eine grosse Menge von Verschlingungen darbietet: es wird dann nothwendig, kürzere Zeitabschnitte jeden für sich besonders zu zeichnen.

Halten wir die Leistungen des neuen Apparats und des Magnetometers zusammen, so ergibt sich, dass beide in Beziehung auf \textit{einige} Zwecke einander wechselseitig ergänzen müssen, in Beziehung auf andere hingegen gleiche Anwendbarkeit haben. Zur Bestimmung der absoluten Declination kann nur das Magneto\SourceLeaf{373}meter dienen, nicht aber der neue Apparat: die Veränderungen der Declination, und besonders die schnell wechselnden lassen sich mit beiden verfolgen. Zur Bestimmung der absoluten Intensität können beide Apparate dienen, obwohl die Anwendung des Magnetometers etwas weniger complicirt ist, als der alleinige Gebrauch des neuen Apparats sein würde; aber jenes für sich allein kann die Intensität nur in ihrem Mittelwerthe während eines gewissen Zeitraumes geben, und die schnell wechselnden Änderungen in demselben entgehen diesem Instrumente gänzlich, während der neue Apparat diese auf das befriedigendste nachweist. Für alle sonstigen Anwendungen, z. B. um Magnetstäbe rücksichtlich ihrer magnetischen Stärke unter einander zu vergleichen; ferner, in Verbindung mit einem Multiplicator, für galvanometrische und telegraphische Zwecke, sind beide gleich brauchbar; ja in den beiden letztern Beziehungen hat der neue Apparat noch einen bedeutenden Vorzug, da man, wie schon bemerkt ist, in seiner Gewalt hat, ihn so nahe man will astatisch zu machen.

Ein paar Proben von der Empfindlichkeit des Apparats als Galvanometer dürfen hier wohl angeführt werden. Der den Magnetstab umgebende Multiplicator enthält 610 Umwindungen mit Seide übersponnenen Kupferdrahts, und ein galvanischer Strom hat in diesem allein schon eine Drahtlänge von mehr als 6000 Fuss zu durchlaufen. Diese Drahtlänge vergrössert sich auf 13000 Fuss, wenn der Strom zugleich nach dem physikalischen Cabinet geht. Gewöhnlich aber werden noch andere Apparate mit in die Kette gebracht, so dass bei manchen Versuchen die ganze Drahtlänge 40000 Fuss oder fast zwei Meilen beträgt. Dabei muss aber noch bemerkt werden, dass bei weitem der grösste Theil dieses Drahts sehr dünner ist, und dass diese Länge, insofern die Stärke des Stroms dadurch bedingt wird, einem etwa 8 Meilen langen Draht von derjenigen Stärke äquivalirt, welche der Verbindungsdraht zwischen der Sternwarte und dem physikalischen Cabinet hat. Trotz der so langen Kette geben nun selbst die schwächsten galvanischen Kräfte dem schweren Magnetstabe eine nicht bloss merkliche, sondern zu scharfen Messungen hinreichende Ablenkung. Dies gilt z. B. vom Thermogalvanismus, in Beziehung auf welchen manche Physiker die irrige Vorstellung haben, als ob er eine sehr lange Kette nicht durchdringen könne. Bei den hiesigen Vorrichtungen, und unter Anwendung eines thermogalvanischen Apparats von eigenthümlicher Construction, reicht die Berührung der Verbindungsstelle mit dem Finger hin, jene Wirkung hervorzubringen.

\SourceLeaf{374}Zu einer andern interessanten Bemerkung gibt die Anwendung auf die gewöhnliche Reibungselektricität Veranlassung. Dass diese, durch einen Multiplicator geleitet, die Magnetnadel auf ganz ähnliche Art ablenkt, wie ein hydrogalvanisch erregter Strom, hat bekanntlich \textsc{Colladon} entdeckt, dessen Anfangs bezweifelte Versuche späterhin durch \textsc{Faraday} bestätigt sind. Der letztere Physiker hat zuerst ins Licht gesetzt, dass in einer sehr starken elektrischen Batterie nicht mehr Elektricität entwickelt ist, als schon sehr geringe hydrogalvanische Erregungsmittel in wenigen Secunden durch einen Leitungsdraht von mässiger Länge treiben. Mit den hiesigen Apparaten war zwar gleichfalls schon vor mehreren Jahren sowohl die Realität, als die geringe Grösse der elektromagnetischen Wirkung der Maschinenelektricität durch Versuche bestätigt gefunden: es schien jedoch der Mühe werth, diese Versuche mit Hülfe des neuen so viel empfindlichern Apparats zu wiederholen. Anstatt eine Leidner Flasche oder eine Batterie von Flaschen durch die Drahtkette zu entladen (wie \textsc{Colladon} und \textsc{Faraday} gethan hatten), wurde nur Conductor und Reibzeug einer im physikalischen Cabinet stehenden Elektrisirmaschine mit den Enden der zur Sternwarte gehenden und mit Inbegriff des Multiplicators 13000 Fuss langen Drahtkette verbunden, und die Elektrisirmaschine anhaltend mit gleichförmiger Geschwindigkeit gedreht; geschah dies mit einer Geschwindigkeit von Einer Umdrehung auf die Secunde, so wurde dadurch der fünfundzwanzigpfündige Magnetstab im neuen Apparat in der Sternwarte in einer Ablenkung von 144 Scalentheilen (etwas über 50 Minuten) erhalten, positiver oder negativer, je nach der Richtung, in welcher die Elektricität den Multiplicator durchströmte, und in den Versuchen zeigte sich alle nur zu wünschende Regelmässigkeit. Aber als besonders merkwürdig erscheint dabei der Umstand, dass die elektromagnetische Wirkung dieselbe blieb, wenn man auch der Kette durch Hineinbringen anderer Apparate eine Länge von einer ganzen Meile gab. Dies könnte ein wesentlicher Unterschied von andern, hydrogalvanisch, thermogalvanisch, oder durch Induction erzeugten Strömen scheinen, deren durch die Grösse der elektromagnetischen Wirkungen sich äussernde Intensität allemal desto kleiner wird, je mehr man die Leitung verlängert. Ich finde aber darin nur eine schlagende Bestätigung der Theorie, welcher zufolge die durch ungleiche elektromagnetische Wirkung sich äussernde ungleiche Intensität zweier galvanischen Ströme nichts weiter ist, als ungleiche Menge in bestimmter Zeit jedem Querschnitt der Leitung durchströmender Elektricität.
\SourceLeaf{375}Bei den andern Erzeugungsarten entwickelt eine gegebene elektromotorische Kraft desto weniger Elektricität in gegebener Zeit, je grösser der Widerstand ist, welchen die längere Kette dem Strome entgegenstellt: bei unserm Versuch hingegen hängt die Menge der bewegten Elektricität bloss von dem Spiel der Maschine ab, und \textit{alle} in Funkenform auf den Conductor überspringende Elektricität muss die ganze Kette, sie mag kurz oder lang sein, durchlaufen, um sich mit der entgegengesetzten des Reibzeugs auszugleichen.

Um auch noch den Vorzug des neuen Apparats vor dem Magnetometer bei der elektromagnetischen Telegraphie nachweisen zu können, wird die Art, \textit{wie} durch galvanische Ströme telegraphische Zeichen hervorgebracht werden, erst etwas näher betrachtet werden müssen.

\expandafter\def\csname HistoricalNoteText375\endcsname{Nach einer mir von Hrn. von \textsc{Humboldt} mitgetheilten Notiz hatte schon zehn Jahre früher \textsc{Bétancourt} eine Drahtkette von Aranjuez nach Madrid gezogen, vermittelst welcher die Entladung einer Leidner Flasche zu einer telegraphischen Signalisirung dienen sollte. Obgleich nähere Umstände über den Erfolg nicht bekannt zu sein scheinen, so ist doch an dem Gelingen eines solchen Versuchs, wenn er zweckmässig ausgeführt wird, nicht zu zweifeln. Aber immer müsste wohl eine solche Methode auf die Signalisirung eines Ja oder Nein auf eine oder ein Paar im Voraus verabredete Fragen beschränkt bleiben.}
Sobald man wusste, dass die Wirkungen einer \textsc{Volta}ischen Säule sich durch eine sehr lange Kette fortpflanzen, lag der Gedanke sehr nahe, diese Naturkräfte zu telegraphischen Zwecken zu benutzen, und schon vor fast 30 Jahren\HistoricalNote{375}, also zu einer Zeit, wo man erst einen kleinen Theil der galvanischen Wirkungen kannte, schlug \textsc{Sömmerring} die Gasentwicklung dazu vor: bei weitem mehr geeignet für zusammengesetzte Signalisirungen sind aber die erst später bekannt gewordenen magnetischen Wirkungen galvanischer Ströme; indessen ist es auffallend, dass seit \textsc{Oersted}'s Entdeckung eine ziemliche Anzahl Jahre verstrichen ist, ehe jemand an diesen Gebrauch gedacht zu haben scheint. Freilich ist ein gründliches Urtheil über die Anwendbarkeit im Grossen nicht möglich ohne eine genaue quantitative Kenntniss der Schwächung galvanischer Ströme in Folge der Länge und Beschaffenheit der Leitungsdrähte, wovon man vor \textsc{Ohm} und \textsc{Fechner} sehr unvollkommene und unrichtige Vorstellungen hatte. Nachdem im Jahr 1833, hauptsächlich um ähnliche Untersuchungen über das Gesetz der Stärke galvanischer Ströme nach Verschiedenheit der Umstände in grossem Maassstabe anstellen zu können, zwischen der hiesigen Sternwarte und dem physikalischen Cabinet eine Drahtverbindung gemacht war, von welcher grossartigen Anlage das Ver\SourceLeaf{376}dienst der sehr schwierigen Ausführung allein dem Herrn Professor \textsc{Weber} gehört, wurde diese Kette gleich vom Anfang an oft zu telegraphischen Zeichen benutzt, nicht bloss zu einfachen, um täglich die Uhren zu vergleichen, sondern versuchsweise auch zu zusammengesetzten; und die Möglichkeit, Buchstaben, Wörter und ganze Phrasen zu signalisiren. wurde dadurch schon damals zu einer evidenten Thatsache\HistoricalNote{376}. Bei diesen Versuchen wurde ein hydrogalvanisch und nur mit schwachen Mitteln, nemlich einem einzigen oder einem doppelten Plattenpaar und ungesäuertem Wasser, erregter Strom angewandt; ich halte mich jedoch nicht dabei auf, das damals gebrauchte Verfahren hier umständlich zu beschreiben, da ich später ein davon ganz verschiedenes an dessen Stelle gesetzt habe. Bei jenem Verfahren blieb die Unbequemlichkeit, dass durch unsere einfache Kette und nach der Einrichtung der Apparate, bei welchen dergleichen Versuche nur eine Nebensache waren, in Einer Minute sich nicht mehr als zwei Buchstaben signalisiren liessen. Auch bei einer abgeänderten bloss für das Telegraphiren berechneten Einrichtung hätte diese Geschwindigkeit (mit welcher übrigens offenbar die Länge der Kette oder die Entfernung der Endpunkte gar nichts zu thun hat) sich nicht viel vergrössern lassen, so lange nur eine einfache Kette angewandt würde, wohl aber in hohem Grade mit einer vielfachen: allein eine solche einzurichten, war hier kein hinlänglicher Beweggrund vorhanden, da theils der Erfolg an sich gar nicht zweifelhaft sein konnte, theils der eigentlich wissenschaftliche Nutzen einer solchen vielfachen Kette mit den bedeutenden Kosten in keinem Verhältniss gestanden haben würde.

Dagegen hat mich die Theorie der Inductionsgesetze auf ein ganz verschiedenes Verfahren geführt, wonach schon seit mehr als zwei Jahren eine einfache Kette mit dem vollkommensten Erfolge zu einem viel schnelleren Telegraphiren dient; und es wird mir um so eher verstattet sein, bei demselben noch etwas zu verweilen, da ich bisher noch nichts Näheres darüber öffentlich bekannt gemacht habe.

Die Vorrichtung, welche ich einen Inductor nenne, habe ich schon vor mehreren Jahren anderwärts beschrieben{\renewcommand{\thefootnote}{**)}\HistoricalNote{376b}}. Ich muss jedoch bemerken, dass anstatt des in der ersten Nachricht beschriebenen Inductors von 1050 Umwindun\SourceLeaf{377}gen, und des nachher auf 3537 Umwindungen verstärkten, gegenwärtig einer von 7000 Umwindungen gebraucht wird, worin die Drahtlänge allein mehr als 7000 Fuss beträgt. Durch eine äusserst einfache Manipulation mit diesem Inductor (dadurch nemlich, dass man ihn von einem doppelten Magnetstab, über welchen er zu Anfang geschoben ist, schnell abzieht und sogleich wieder, ohne ihn umzukehren in die vorige Lage zurückbringt) wird bewirkt, dass schnell nach einander zwei starke entgegengesetzte galvanische Ströme durch den Leitungsdraht gehen, deren jeder nur eine äusserst kurze Zeit dauert. Die Wirkung dieser beiden Ströme auf eine wo immer in der Kette befindliche, von einem Multiplicator umgebene Magnetnadel besteht darin, dass dieser für einen Augenblick eine sehr lebhafte Geschwindigkeit ertheilt, aber dann sogleich vollkommen wieder aufgehoben wird. Die Nadel macht also eine sehr lebhafte, aber nur kleine Bewegung, nach Gefallen rechts oder links, und steht dann sogleich wieder ganz still.

Dass sich nun die Abwechslungen solcher zuckenden Bewegungen auf mancherlei Art combiniren und zur Signalisirung von Buchstaben benutzen lassen, ist von selbst klar. Die Zeichen möglichst schnell und präcis zu geben, so wie, von der andern Seite, sie mit Leichtigkeit und Schnelligkeit zu lesen, wird allerdings eine gewisse Einübung erforderlich sein: aber auch schon, ohne sich eine solche besonders angeeignet zu haben, kann man, wie öftere Erfahrungen gezeigt haben, in Einer Minute füglich etwa sieben Buchstaben signalisiren. Wollte man für die Manipulation eigne \textit{mechanische} Vorrichtungen treffen, so würde sich ohne Zweifel die Schnelligkeit und Präcision noch bedeutend erhöhen lassen.

Gerade bei dieser Art des Telegraphirens hat nun der neue Apparat einen bedeutenden Vorzug vor dem Magnetometer, und zwar wegen folgender Umstände. Obgleich die beiden entgegengesetzten Impulse, aus welchen Ein einfaches Zeichen besteht, ihrer \textit{Stärke} nach \textit{genau gleich} sind, und daher der zweite genau eben so viel Geschwindigkeit vernichtet, als der erste hervorgebracht hat, so kann dennoch die Nadel zwischen den Zeichen nicht in absoluter Ruhe sein, weil diese nur da möglich ist, wo jene sich in ihrer natürlichen Gleichgewichtsstellung befindet. Ist sie auch, \textit{vor} einem Zeichen, in dieser Stellung, so wird sie doch eben \textit{durch} das Zeichen etwas, wenn auch nur wenig, daraus verrückt, und die auf die Nadel wirkende Directionskraft strebt dann, sie nach derselben zurückzuführen. Wenn nun gleich so, in Folge Eines Zeichens, nur eine äusserst
\SourceLeaf{378}schwache Bewegung entstehen kann, so wird doch nach einer grossen Menge von Zeichen durch Anhäufung eine beträchtliche Entfernung von der natürlichen Gleichgewichtsstellung eintreten können, mithin in Folge derselben auch zwischen den Zeichen so viel Bewegung, dass die Zeichen dadurch etwas von ihrer scharfen Ausprägung verlieren. Diese Störung tritt nun, wie man bei einiger Überlegung leicht einsieht, unter sonst gleichen Umständen nachtheiliger hervor, wenn die Nadel, an deren zuckenden Bewegungen die Zeichen beobachtet werden, eine kurze, als wenn sie eine lange Schwingungsdauer hat, daher mehr an dem Magnetometer des magnetischen Observatoriums, als an dem in der Sternwarte aufgehängt gewesenen mit fünfundzwanzigpfündiger Nadel; noch weniger hingegen, als bei letzterem, an dem neuen jetzt dessen Stelle einnehmenden Apparat, wenn dessen Magnetstab in der verkehrten Lage zu einer fast astatischen Nadel eingerichtet ist. In der That wird dann dieselbe, selbst nach einer beträchtlichen Entfernung von ihrer Gleichgewichtsstellung, von der sie dahin zurücktreibenden, vergleichungsweise schwachen, Directionskraft in keine die Zeichen erheblich störende Bewegungen versetzt, während der Strom im Multiplicator eben so stark auf sie wirkt, und also eben so grosse Zuckungen hervorbringt, als gehörte sie zu einem gewöhnlichen Magnetometer.

Gegen die Nachtheile und Unbequemlichkeiten unzeitiger Schwingungsbewegungen, sowohl bei dieser Art des Telegraphirens, als bei manchen andern Anwendungen der magnetischen Apparate, leistet übrigens eine eigne Vorrichtung, die ich vor kurzem habe ausführen lassen, ungemein nützliche Dienste. Ich nenne diese Vorrichtung einen \textit{Dämpfer}, da ihre Wirkung darin besteht, Schwingungsbewegungen, die sonst mit sehr langsamer Abnahme viele Stunden fortdauern würden, in sehr kurzer Zeit ganz zu vernichten. Diese Wirkung leistet der vorerst nur für das Magnetometer des magnetischen Observatoriums angefertigte Dämpfer in ganz eminentem Grade, so dass die grössten Schwingungsbewegungen in wenigen Minuten gänzlich erlöschen. Eine ähnliche Vorrichtung kann aber bei jeder schwingenden Nadel, bei einem Magnetometer oder bei dem neuen hier in Rede stehenden Apparate angebracht werden, und wird bei allen Apparaten, die zum Telegraphiren nach der hier beschriebenen Methode ernstlich angewandt werden sollen, einen wesentlichen Bestandtheil ausmachen müssen. Eine ausführlichere Erklärung dieser Vorrichtung würde aber von dem gegenwärtigen Gegenstande zu weit abführen.

\SourceLeaf{379}In obigem ist dem neuen Apparate noch keine besondere Benennung beigelegt. Nach seiner wichtigsten Anwendung könnte man ihn einen Intensitätsmesser nennen. In so fern er aber zu eben so mannichfaltigen scharfen magnetischen Messungen dient, wie das Magnetometer, hätte er wohl eben so gut auf \textit{dieselbe} Benennung Anspruch. Der wesentlichste Unterschied ist der, dass der neue Apparat an \textit{zwei} Fäden aufgehängt ist, wodurch eben eine \textit{neue} Directionskraft gewonnen wird, mit welcher die magnetische commensurabel ist. Die übrigen Unterschiede, namentlich die Art, wie der Spiegel angebracht ist, ferner die Mittel zur Messung der Drehung der einzelnen Bestandtheile gegen einander, sind nothwendige sich von selbst ergebende Bedingungen für die zu erreichenden Zwecke. Man könnte daher dem neuen Apparat ein \textit{Bifilar-} oder \textit{Bipensil-Magnetometer} nennen, um es von dem ältern Instrumente, dem einfachen oder Unifilar-Magnetometer zu unterscheiden. Ich darf wohl meine Überzeugung aussprechen, dass einer allgemeinern Verbreitung desselben, und besonders einer Anwendung in den Terminsbeobachtungen neben dem einfachen Magnetometer an mehrern weit von einander entlegenen Orten, bedeutende Fortschritte unsrer Kenntniss der wunderbaren Störungen des Erdmagnetismus bald folgen werden.

\clearpage
\fancyhead[L]{\small SCHWINGUNGSDAUER EINER MAGNETNADEL.}
\SourceLeaf{380}
\pdfbookmark[0]{Anleitung zur Bestimmung der Schwingungsdauer einer Magnetnadel. 1837}{schwingungsdauer-1837}
\begin{center}
{\large ANLEITUNG ZUR BESTIMMUNG DER SCHWINGUNGSDAUER EINER MAGNETNADEL.}
\par\medskip
Resultate aus den Beobachtungen des magnetischen Vereins. 1837. IV.
\end{center}

Die Aufgabe, zu deren Auflösung hier eine Anleitung gegeben werden soll, hat ein mehrseitiges Interesse. Eine wenn auch noch nicht sehr genaue Kenntniss der Schwingungsdauer ist schon zur Ausübung der für die Bestimmung des Ruhestandes der Nadel gegebenen Vorschriften nothwendig (\textit{Res.} von 1836, II.): zur Ausmittelung der absoluten Intensität des Erdmagnetismus hingegen ist der auf das schärfste bestimmte Werth der Schwingungsdauer ein wesentliches Element. Aber auch an sich kann die Ausübung der zur Bestimmung der Schwingungsdauer gehörenden Operationen wie eine gute Vorübungsschule für astronomische Beobachtungen betrachtet werden, da jene namentlich mit den Beobachtungen der Sterndurchgänge am Mittagsfernrohr die grösste Ähnlichkeit, aber vor denselben den Vorzug haben, dass sie grösserer Schärfe fähig sind, und, durch ungünstigen Luftzustand ungestört, jede Stunde nach Gefallen vorgenommen werden können. Es scheinen daher auch solche Beobachtungen am Magnetometer besonders dazu geeignet zu sein, über einen bisher noch nicht genügend aufgeklärten Gegenstand Licht zu verbreiten, nemlich über die constanten Differenzen zwischen den Resultaten verschiedener Beobachter am Mittagsfernrohr, welche doch nur daher rühren können, dass die optischen oder die akustischen Eindrücke oder beide, bei verschiedenen Personen und nach Verschiedenheit der Umstände nicht gleichzeitig ins Bewusstsein kommen.

Die Schwingungsdauer einer Nadel ist die Zwischenzeit zwischen zwei auf einander folgenden äussersten Stellungen (Elongationen) derselben. Die Nadel
\SourceLeaf{381}befindet sich in jeder Elongation streng genommen ohne alles Verweilen; allein, da die Geschwindigkeit der Bewegung bis zum Verschwinden nach der Stetigkeit abnimmt, und eben so von da an wieder zunimmt, so erscheint sie für unsere Sinne um die Zeit der Elongation mit einer grössern oder geringern Dauer als ruhend, welche aber freilich als solche bei kurzen Schwingungszeiten und grossen Bögen kaum erkannt wird. In allen Fällen aber bleibt eine solche unmittelbare Auffassung des Zeitpunkts der Elongation an Schärfe weit zurück gegen eine mittelbare Bestimmung durch correspondirende Beobachtungen, indem man nemlich das Mittel aus den beiden Zeiten nimmt, wo ein und derselbe Theilstrich der Scale beim Hin- und Rückgange auf dem Vertikalfaden des Fernrohrs erscheint.

Im Allgemeinen ist es am vortheilhaftesten, dazu einen Theilstrich in oder nahe bei der Mitte des Schwingungsbogens zu wählen, theils weil da die Bewegung am schnellsten, mithin die Beobachtung der Zeit selbst am schärfsten ist, theils weil beim Beobachten mehrerer auf einander folgender Schwingungen die Zwischenzeiten zwischen den Aufzeichnungen nahe gleich werden. In einzelnen Fällen, namentlich bei sehr langer Schwingungsdauer, kann es übrigens allerdings zuweilen vortheilhaft sein, andere oder selbst mehrere verschiedene Scalenstellen anzuwenden, was jedoch hier bei Seite gesetzt bleiben kann.

Bei kleinen Schwingungen thut man wohl, den der Mitte nächsten Theilungspunkt selbst zu wählen, bei grössern ziehe man den bequemer zu beachtenden nächsten Theilstrich bei den Fünfern oder Zehnern der Scale vor; bei sehr grossen Schwingungen hingegen wird es wegen der grossen Schnelligkeit, mit welcher die Mitte der Scale durch das Gesichtsfeld geht, nothwendig, die gewählte Stelle der Scale, etwa durch einen nicht zu feinen über die Scale gehängten schwarzen Faden, gehörig augenfällig zu machen.

Wenn der Vorübergang am Faden nicht genau mit einem Secundenschlage zusammenfällt, so setzt man den Bruchtheil nach dem geschätzten Verhältniss der beiden Entfernungen an, in welchen die betreffende Scalenstelle vom Faden beim vorhergehenden und folgenden Sekundenschlage erscheint, eben so wie es die meisten Astronomen beim Beobachten am Mittagsfernrohr gewohnt sind. Man theilt also, unmittelbar, nicht die Zeit, sondern den Raum.

Die Bestimmung der Schwingungsdauer aus einer einzigen Schwingung ist natürlich nur einer sehr beschränkten Schärfe fähig; man gründet deshalb jene immer auf eine grössere Anzahl auf einander folgender Schwingungen. Zwar ist
